% Part: sets-functions-relations
% Chapter: functions
% Section: partial-functions
% Hindi translation (hi-Deva-IN) of Open Logic commit 9620cc73f9c8e0ad003c514a5d3748f29611c4c0.

\documentclass[../../../include/open-logic-section]{subfiles}


\begin{document}

\olfileid[hi]{sfr}{fun}{par}

\olsection{आंशिक फलन}

\begin{explain}
कभी-कभी फलन की परिभाषा को इस प्रकार शिथिल करना उपयोगी होता है कि फलन का
निर्गत सभी संभावित निवेशों के लिए परिभाषित होना आवश्यक न हो। ऐसे
प्रतिचित्रणों को \emph{आंशिक फलन} कहते हैं।
\end{explain}

\begin{defn}
\emph{आंशिक फलन} $f \colon A \pto B$ ऐसा प्रतिचित्रण है जो~$A$ के प्रत्येक
!!{element} को~$B$ का अधिक-से-अधिक एक !!{element} देता है। यदि $f$,~$B$ का
कोई अवयव $x \in A$ को देता है, तो हम कहते हैं कि $f(x)$ \emph{परिभाषित} है;
अन्यथा वह \emph{अपरिभाषित} है। यदि $f(x)$ परिभाषित है, तो हम
$f(x) \fdefined$ लिखते हैं; अन्यथा $f(x) \fundefined$। आंशिक फलन~$f$ का
\emph{प्रांत},~$A$ का वह उपसमुच्चय है जहाँ वह परिभाषित है; अर्थात
$\dom{f} = \Setabs{x \in A}{f(x) \fdefined}$।
\end{defn}

\begin{ex}
प्रत्येक फलन $f\colon A \to B$ एक आंशिक फलन भी है। जो आंशिक फलन~$A$ पर
हर जगह परिभाषित हैं---अर्थात जिन्हें अब तक हम बस फलन कहते आए हैं---उन्हें
\emph{पूर्ण} फलन भी कहते हैं।
\end{ex}

\begin{ex}
$f \colon \Real \pto \Real$ द्वारा दिया गया आंशिक फलन $f(x) = 1/x$,
$x = 0$ के लिए अपरिभाषित और अन्य सभी स्थानों पर परिभाषित है।
\end{ex}

\begin{prob}
$f\colon A \pto B$ दिया हो, तो आंशिक फलन $g\colon B \pto
A$ को इस प्रकार परिभाषित कीजिए: किसी भी $y \in B$ के लिए, यदि कोई अद्वितीय
$x \in A$ ऐसा है कि $f(x) = y$, तो $g(y) = x$; अन्यथा
$g(y) \fundefined$। दिखाइए कि यदि $f$ एकैकी है, तो सभी
$g(f(x)) = x$, सभी $x \in \dom{f}$ के लिए, और
$f(g(y)) = y$, सभी $y \in \ran{f}$ के लिए।
\end{prob}

\begin{defn}[आंशिक फलन का आलेख]
मान लीजिए कि $f\colon A \pto B$ एक आंशिक फलन है।~$f$ का \emph{आलेख},
निम्न प्रकार परिभाषित संबंध $R_f \subseteq A \times B$ है:
\[
R_f = \Setabs{\tuple{x,y}}{f(x) = y}.
\]
\end{defn}

\begin{prop}
मान लीजिए कि $R \subseteq A \times B$ में यह गुण है कि जब भी $Rxy$
और $Rxy'$ हों, तब $y = y'$। तब $R$ उस आंशिक फलन $f\colon A \pto B$ का
आलेख है जिसे इस प्रकार परिभाषित किया जाता है: यदि कोई $y$ ऐसा है कि $Rxy$,
तो $f(x) = y$; अन्यथा $f(x) \fundefined$। यदि $R$ \emph{क्रमिक} भी है,
अर्थात प्रत्येक $x \in A$ के लिए कोई $y \in B$ ऐसा है कि $Rxy$, तो $f$
पूर्ण है।
\end{prop}

\begin{proof}
मान लीजिए कि कोई $y$ ऐसा है कि $Rxy$। यदि कोई अन्य $y'
\neq y$ ऐसा होता कि $Rxy'$, तो $R$ पर लगी शर्त का उल्लंघन होता। अतः यदि
कोई $y$ ऐसा है कि $Rxy$, तो वह $y$ अद्वितीय है; इसलिए $f$ सुपरिभाषित है।
स्पष्टतः, $R_f = R$ और $f$ पूर्ण है यदि~$R$ क्रमिक है।
\end{proof}

\end{document}
